\scanpage{130}
% archive sheet number: 129
aux op. $\sigma$, on trouve que l'\deleted{app}image de $Y$ dans $E$ \deleted{est} (étant \unsure{invariant|invariante} sous $V$ et non vide) est $V$ tt entier. Or $\inserted{\operatorname{card}Y=}4c_\nu=4\deleted{\unsure{2^\nu=2^2}}$ \deleted{\unsure{=}}\inserted{et} $\operatorname{card}E=\operatorname{card}V$ \deleted{\unsure{donc}} $\inserted{=2^{c_\nu+1}}=2^2=4$, donc \deleted{\unsure{pour avoir un}} $\varphi$ est bijectif.

b) Si $\nu\ne0$ i.e. $\nu=1$ ou $2$, alors \deleted{l'image} $u$ a comme noyau \deleted{Im} $\delta(\mathbb F_2)=\{0,z\}$, comme image $\operatorname{Ker}\unsure{\eta|\gamma}=V$ \deleted{\unsure{A dire vrai} \inserted{\unsure{(facile)|(faible)}} \unsure{puis} $Y$ est un \unsure{tors.}}
